\documentclass[11pt]{article}
\usepackage{icmlko}
\begin{document}

\title{거래가 없는 채택 --- 유보를 안 보고 고른 가림 한 칸이 자 다섯을 모두 올린다}
\author{월드모델 실험실}
\date{노트 498 · 2026년 8월 3일}
\maketitle

\begin{abstract}
노트 497 이 ``버는 곳만 가리면 $+0.0161$''이라는 상한을 냈는데 \emph{어느
곳이 버는지를 유보로 골랐다}. 이번에는 \textbf{유보를 안 보고} 골랐다 ---
학습 안에서 시간을 한 번 더 잘라(안쪽 학습 $<\!2023$ · 안쪽 검증
$2023\!\sim\!2025$) 후보 축 열 개를 하나씩 전 도메인에서 가려 (도메인, 축)
차를 재고, 문턱을 넘는 짝만 남긴다.

\textbf{문턱 하나로 답이 뒤집힌다.}

\begin{center}
\begin{tabular}{llrr}
\hline
문턱 & 뽑힌 짝 & 판 차 & 부호 \\
\hline
차 $> 0.005$ & 아홉 & $-0.0085$ & $0/12$ \\
차 $\geq 0.010$ · $n_{\text{val}} \geq 150$ & \textbf{하나}(애니 트렌드) & $\mathbf{+0.0046}$ & $\mathbf{12/12}$ \\
\hline
\end{tabular}
\end{center}

그리고 그 하나에 자 다섯을 다 걸었더니 \textbf{전부 양수다}(짝 씨앗 $10$).

\begin{center}
\begin{tabular}{lrrl}
\hline
자 & 차 & 부호 & $95\%$ 구간 \\
\hline
판 & $+0.0045$ & $10/10$ & $[+0.0033, +0.0056]$ \\
날짜 통제 판 & $+0.0045$ & $10/10$ & $[+0.0036, +0.0054]$ \\
무리 안 판 & $\mathbf{+0.0088}$ & $10/10$ & $[+0.0078, +0.0097]$ \\
KR 만화(집 밖) & $\mathbf{+0.0152}$ & $10/10$ & $[+0.0125, +0.0179]$ \\
비게임 앱(집 밖) & $+0.0039$ & $8/10$ & $[+0.0016, +0.0062]$ \\
\hline
\end{tabular}
\end{center}

\textbf{거래가 없다.} 노트 426 의 거울(판을 잃고 집 밖을 사는 규약)을 쓸
일이 없는 첫 채택이다. \textbf{채택한다} --- \texttt{F18\_bagboost.DOMDROP
= \{"애니": ("trend\_",)\}}.
\end{abstract}

\begin{figure}[t]
\centering
\includegraphics[width=\textwidth]{figs/f1.pdf}
\caption{\textbf{가} --- 두 문턱. \textbf{나} --- 자 다섯.
\textbf{다} --- 도메인별.}
\end{figure}

\section{왜 아홉 짝은 졌나}

$1$ 차에서 뽑힌 아홉 중 \textbf{맞은 것은 둘뿐}이다(애니 트렌드 $+0.0334$
· 시장팝업 $+0.0015$). 나머지는 아이돌 $-0.0877$ · 게임 $-0.0606$ ·
모바일 $-0.0593$ · 팝업 $-0.0203$ 로 크게 틀렸다.

크기만의 문제는 아니다 --- 안쪽 검증이 제일 큰 웹툰($851$ 행)도 틀렸다.
그러나 제일 작은 둘(팝업 $24$ · 아이돌 $45$)이 제일 크게 틀렸고, 그 둘의
안쪽 차가 $+0.07 \cdot +0.05$ 로 \emph{제일 커 보였다.} \textbf{작은
표본이 큰 효과를 흉내 낸다.}

노이즈를 감안한 자연스러운 통계 --- 안쪽 차 $\times \sqrt{n_{\text{val}}}$
--- 로 줄을 세우면 \textbf{애니 트렌드가 $0.435$ 로 $1$ 등}이고 팝업
venue($0.357$) · 아이돌 위키($0.315$)가 뒤따른다. \emph{이 통계는 유보를
전혀 안 본다.} 다만 우리는 그것을 사후에 계산했으므로 \textbf{다음
사이클의 제안으로만 적는다.}

\begin{quote}
\textbf{정직하게} --- 채택에 쓴 $n_{\text{val}} \geq 150$ 은 $1$ 차가
어디서 틀렸는지 \emph{본 뒤에} 골랐다. 그러니 이 채택의 근거는 ``자동
규칙이 작동한다''가 아니라 \textbf{``학습만 보고도 제일 큰 짝은 짚힌다''}
이다. 규칙 자체는 아직 검증 안 됐다.
\end{quote}

\section{가림은 도메인 안에 안 갇힌다}

애니 하나만 건드렸는데 열한 도메인이 다 움직인다 --- 도서 $+0.0242$ ·
아이돌 $-0.0577$ · 팝업 $-0.0253$. 공유 열이 달라지기 때문이다. 애니가
판 무게 $17.7\%$ 이고 아이돌 $\cdot$ 팝업이 $1.5\% \cdot 1.9\%$ 라 순합이
양수다.

\begin{quote}
\textbf{조항} --- \textbf{도메인별 가림의 값은 그 도메인에서 재면 안
된다.} 공유 축을 쓰는 모형에서는 한 칸을 바꾸면 모든 칸이 바뀐다. 이것이
$1$ 차 고르개가 진 진짜 이유이기도 하다 --- 고르개는 자기 도메인 효과만
보고 이웃 값을 못 본다.
\end{quote}

\section*{배선}

\texttt{F18\_bagboost.DOMDROP} + \texttt{\_domdrop()} 를 \texttt{fit} 과
\texttt{predict} 양쪽에 넣었다. 가드 --- 애니 트렌드 값을 난수로 바꿔도
예보 최대 차 $0.00\mathrm{e}{+}00$, 웹툰(대조)은 $5.3\mathrm{e}{+}02$.
판 $0.4812 \to \mathbf{0.4857}$. 되돌림은 ``\texttt{DOMDROP} 이 빈 F18''.

\section*{남은 것}

\begin{itemize}
\item \textbf{$\sqrt{n}$ 가중 고르개} --- 다음 사이클에 사전 등록하고
돌린다. 통과하면 규칙이 되고, 그러면 남은 여덟 축까지 열린다.
\item \textbf{이웃 값을 보는 고르개} --- 안쪽 검증에서 \emph{판}을 보고
고르면 이웃까지 값이 매겨진다. $1$ 차는 도메인별로만 봤다.
\item \textbf{왜 애니가 트렌드를 싫어하나} --- 네 번째 실패. 이제
\emph{몰라도 쓰고 있다.}
\end{itemize}

\end{document}
